\subsection{P5: miejsca wyboru w aktach prawnych}
\label{sec:rescorpus}

Tabela \ref{tab:corpus} i rysunek \ref{fig:corpus} podają wynik sprawdzarki na korpusie.
Akty prawne mają od \pkCorpusActMin{} do \pkCorpusActMax{} trafień na 1000 słów włącznie.
Teksty pisane w PK mają od \pkCorpusPkMin{} do \pkCorpusPkMax{} trafień na 1000 słów
włącznie. Korpus aktów liczy \pkCorpusActWords{} słów.

W aktach prawnych przeważają 4 kontrole. Kontrola A3 zgłasza zaimki. Kontrola KW zgłasza
słowa \pkq{może} i \pkq{musi} poza zapisem słów kluczowych. Kontrola B8 zgłasza zdania
dłuższe niż 25 słów. Kontrola A5 zgłasza słowa nieostre. Udział zdań dłuższych niż 25 słów wynosi w aktach
od \pkCorpusOverMin{} \% do \pkCorpusOverMax{} \% włącznie. Podział aktu na zdania jest
przybliżony, bo akt składa się z jednostek redakcyjnych, a nie ze zdań.

Ręczna ocena próby potwierdza trafność wskazań: \pkPrecValid{} z \pkPrecN{} trafień
wskazuje konstrukcję objętą regułą. Precyzja wynosi \pkPrec{} \%, a przedział Wilsona 95 \%
to od \pkPrecLo{} \% do \pkPrecHi{} \% włącznie. Czułości sprawdzarki praca nie mierzy.

Wynik wymaga 3 zastrzeżeń. Po pierwsze, trafienie wskazuje miejsce wyboru, a nie błąd
aktu. Język prawny ma własne konwencje i własne procedury wykładni. Po drugie, nieostrość
tekstu prawnego bywa celowa: słowo \pkq{niezwłocznie} zostawia sądowi ocenę przypadku
(sekcja \ref{sec:limits}). Po trzecie, teksty pisane w PK także mają trafienia. Standard
PK 0.1 zawiera \pkCorpusStdCfive{} zdania z 2 przeczeniami, czyli narusza własną regułę C5.
Część trafień w tekstach PK jest fałszywa: sprawdzarka zgłasza przymiotnik w nazwie
\pkq{małe litery} jako słowo nieostre.

\begin{table}[tp]
\centering
\caption{P5: wynik sprawdzarki na korpusie. Mediana: mediana długości zdania w słowach.
Kolumna 4: udział zdań dłuższych niż 25 słów w procentach. Trafienia: trafienia
\pkLintChecks{} kontroli na 1000 słów. Kolumna końcowa: liczba wystąpień słowa \pkq{się}
na 1000 słów; słowo \pkq{się} jest kandydatem do oceny, a nie naruszeniem.}
\label{tab:corpus}
\small
\begin{tabular}{@{}lrrrrr@{}}
\toprule
Tekst & słowa & mediana & ponad 25 (\%) & trafienia & \pkq{się}\\
\midrule
\wiersze{data/table_corpus.tex}
\bottomrule
\end{tabular}
\end{table}

\rysunek{corpus}{Lingwistyka korpusowa: liczba trafień sprawdzarki na 1000 słów w aktach
prawnych i w tekstach pisanych w PK.}{fig:corpus}
